\documentclass[aspectratio=1610]{beamer}
\usepackage[utf8]{inputenc}
\usepackage{amsmath, amssymb}

\title{Deljenje tajni}
\author{Staša Đorđević}
\date{\today}

\usetheme[color=nudarkblue,numbering=progressbar]{wildcat}
 
\begin{document}

\begin{frame}
\titlepage
\end{frame}
\begin{frame}{Sadržaj}
\tableofcontents
\end{frame}
 
\section{Uvod i motivacija}
 
\begin{frame}{Motivacioni primer - vojni štab}
\begin{block}{Scenario}
Generalni štab čuva šifru za aktivaciju vojnog stanja.\\[0.4em]
Učesnici: \textbf{M} (Ministar), \textbf{N} (Načelnik),
\textbf{K} (Komandant), \textbf{O} (Obaveštajni direktor).
\end{block}
 
\vspace{0.6em}
\textbf{Pravila:}
\begin{enumerate}
\item $O$ \emph{ne sme sam} da aktivira vojnog stanje.
\item $M$ i $O$ \textbf{zajedno mogu}.
\item $N$, $K$ i $O$ \textbf{zajedno mogu}.
\end{enumerate}
 
\vspace{0.8em}
\centering
\large
Minimalni ovlašćeni skupovi: \quad
$\{M,\,O\}$ \quad i \quad $\{N,\,K,\,O\}$
\end{frame}

\begin{frame}{Šta je deljenje tajni?}
\textbf{Problem:} kako podeliti tajnu vrednost $s$ između $n$ učesnika tako da:
\begin{itemize}
\item samo \textbf{određene grupe} mogu rekonstruisati tajnu,
\item sve ostale grupe \textbf{ne saznaju ništa} o njoj.
\end{itemize}
 
\end{frame}
 
\begin{frame}{Osnovni pojmovi}
\begin{itemize}
\item \textbf{Udeo} (share) --- deo tajne dodeljen učesniku
\item \textbf{Kvalifikujući skup} --- skup koji može rekonstruisati tajnu
\item \textbf{Struktura pristupa} $\Gamma$ --- skup svih kvalifikujućih skupova
\item \textbf{Minimalni kvalifikujući skupovi} $m(\Gamma)$
\end{itemize}
 
\vspace{1em}
Za naš primer:
\[m(\Gamma) = \bigl\{\{M,O\},\, \{N,K,O\}\bigr\}\]
\[\Gamma = \bigl\{\{M,O\}, \{N,K,O\}, \{M,O,K\}, \{M,O,N\}, \{M,N,K,O\}\bigr\}\]
\end{frame}

\begin{frame}{Monotona struktura}
\begin{block}{Definicija}
Monotona struktura na skupu $P$ je kolekcija $\Gamma$ podskupova od $P$ takva da:
\begin{itemize}
\item $P \in \Gamma$
\item ako $A \in \Gamma$ i $A \subseteq B \subseteq P$, onda $B \in \Gamma$
\end{itemize}
\end{block}
 
\textbf{Intuicija:} ako grupa može rekonstruisati tajnu, svaka veća grupa to svakako može.
 
\vspace{1em}
 
Šema deljenja tajni = par algoritama $\mathbf{Share}$ i $\mathbf{Recombine}$.
\end{frame}
 
\begin{frame}{Bezbednost}
 
\vspace{1em}
\textbf{Informaciono-teorijska sigurnost} - ni uz neograničenu računarsku moć, neovlašćeni skup ne saznaje ništa.
\end{frame}


\begin{frame}{Prahovne (threshold) strukture}
Poseban i najvažniji slučaj monotone strukture:
\[\Gamma_{t,n} = \{ B \subseteq P : |B| \geq t \}\]
 
 
\begin{itemize}
\item nazivamo je \textbf{$t$-od-$n$ šema}
\item Primer: 2-od-4 šema --- bilo koja dva od četiri učesnika
\end{itemize}
 
\vspace{1em}
\alert{Napomena:} vojni štab \textbf{nije} pragovna struktura ($\{M,N\}$ ne mogu, a $\{N,K,O\}$ mogu).
\end{frame}

\section{Opšte šeme}
 
\begin{frame}{Dve opšte konstrukcije}
Rade za \textbf{bilo koju} monotonu strukturu pristupa, ali su neefikasne.
\vspace{1em}
\begin{enumerate}
\item Ito-Nišizeki-Saito (INS) šema
\item Replicirano deljenje tajni
\end{enumerate}
\vspace{1em}
\end{frame}

 

\begin{frame}{INS šema - DNF formula}
Koristi disjunktivnu normalnu formu strukture pristupa:
\[(M \wedge O) \vee (N \wedge K \wedge O)\]
 
\textbf{Ideja:} svako "ili" $\to$ nezavisno deljenje, svako "i" $\to$ XOR dekompozicija.
\[s = s_1 \oplus s_2 \qquad s = s_3 \oplus s_4 \oplus s_5\]
\end{frame}
 
\begin{frame}{INS šema - udeli}
\[s_M = s_1, \quad s_N = s_3, \quad s_K = s_4, \quad s_O = s_2 \| s_5\]
 
\begin{itemize}
\item $\{M, O\}$: imaju $s_1, s_2 \Rightarrow s_1 \oplus s_2 = s$ \checkmark
\item $\{N,K,O\}$: imaju $s_3,s_4,s_5 \Rightarrow s_3 \oplus s_4 \oplus s_5 = s$ \checkmark
\item $\{M, N\}$: imaju $s_1, s_3$ --- ne mogu! \checkmark
\end{itemize}
 
\vspace{1em}
\alert{Problem:} $O$ čuva \textbf{dva} udela --- duplo više podataka.
\end{frame}



\begin{frame}{Replicirano deljenje tajni}
\textbf{Ideja:} koristi \emph{maksimalne nekvalifikujuće} skupove $A_i$ i njihove komplemente $B_i = P \setminus A_i$.
 
\[A_1 = \{M,N,K\} \Rightarrow B_1 = \{O\}\]
\[A_2 = \{N,O\} \Rightarrow B_2 = \{M,K\}\]
\[A_3 = \{K,O\} \Rightarrow B_3 = \{M,N\}\]
 
\[s_M = s_2 \| s_3, \quad s_N = s_3, \quad s_K = s_2, \quad s_O = s_1\]
\end{frame}
 

\begin{frame}{Problem efikasnosti}
 
\vspace{1em}
Obe šeme su informaciono-teorijski sigurne, ali veličina udela
raste \textbf{eksponencijalno} sa brojem učesnika.
 
\vspace{0.8em}
\alert{Potrebna nam je efikasnija šema za pragovne strukture.}
$\to$ Šamirova šema
\end{frame}


\section{Rid-Solomonovi kodovi}
 
\begin{frame}{Rid-Solomonovi kodovi}
Radimo nad konačnim poljem $\mathbb{F}_q$, biramo $n$, $t$,
i tačke $X = \{x_1, \ldots, x_n\} \subset \mathbb{F}_q \setminus \{0\}$.
 
\vspace{0.6em}
Posmatramo polinome stepena $\leq t$:
\[f(X) = f_0 + f_1 X + \cdots + f_t X^t, \qquad f_i \in \mathbb{F}_q.\]
 
\textbf{Kodna reč:}
\[c = \bigl(f(x_1),\, f(x_2),\, \ldots,\, f(x_n)\bigr)\]
 
\vspace{0.4em}
\textbf{Primer:} $q=101$, $n=7$, $t=2$, $f(X) = 20 + 57X + 68X^2$:
\[c = (44,\; 2,\; 96,\; 23,\; 86,\; 83,\; 14)\]
\end{frame}
 
\begin{frame}{Lagranžova interpolacija}
Polinom stepena $\leq t$ jednoznačno je određen vrednošću u $t+1$ tačaka.
 
\vspace{0.8em}
\textbf{Lagranžovi bazni polinomi:}
\[\delta_i(X) = \prod_{j \neq i} \frac{X - x_j}{x_i - x_j}
\qquad \Rightarrow \qquad
\delta_i(x_i) = 1, \quad \delta_i(x_j) = 0 \text{ za } j \neq i\]
 
\textbf{Rekonstrukcija:}
\[f(X) = \sum_{i=1}^{n} c_i \cdot \delta_i(X)\]
 
\end{frame}

 \begin{frame}{Detekcija greške}
\textbf{Ideja:} ako kodna reč sadrži grešku, Lagranžova interpolacija
daje polinom stepena većeg od $t$.
 
\vspace{0.8em}
\textbf{Primer:} primamo $c' = (44, 2, \mathbf{25}, 23, 86, 83, 14)$
umesto ispravnog $c = (44, 2, 96, 23, 86, 83, 14)$.
 
\vspace{0.6em}
Interpolacija nad svih 7 tačaka daje polinom stepena 6, a ne $t = 2$.
 
\vspace{0.6em}
$\Rightarrow$ Zaključujemo da postoji greška, ali ne znamo gde.
\end{frame}

\begin{frame}{Berlekamp-Velč algoritam}
\textbf{Ideja:} nelinearan problem (ne znamo ni $f$ ni koje tačke su pogrešne)
pretvaramo u linearan sistem.
 
\vspace{0.6em}
Tražimo $Q(X,Y) = f_0(X) - f_1(X) \cdot Y$ uz uslove
$\deg f_0 \leq 2t$, $\deg f_1 \leq t$, $f_1(0) = 1$.
 
\vspace{0.6em}
\textbf{Korak 1:} nametnemo $Q(x_i, y_i) = 0$ za sve primljene tačke
$\to$ linearni sistem po koeficijentima $f_0$ i $f_1$.
 
\textbf{Korak 2:} rešimo sistem i izračunamo:
\[f(X) = \frac{f_0(X)}{f_1(X)}\]
 
\textbf{Zašto važi:} polinom $P(X) = f_0(X) - f_1(X) \cdot f(X)$
ima stepen $\leq 2t$, ali se poništava u svim ispravnim tačkama
(kojih je više od $2t$), pa mora biti identički nula.
\end{frame}



\section{Šamirova šema}
 
\begin{frame}{Šamirova šema - ideja}
\begin{block}{Motivacija (Šamir, 1979)}
11 naučnika želi da zaključa sef koji se otvara samo ako je
prisutno bar 6 naučnika. Naivno rešenje (brave i ključevi) nije efikasno.
\end{block}
 
\vspace{0.8em}
Rešenje: polinom stepena $t$ jednoznačno određen sa $t+1$ tačaka.
\[\text{polinom stepena } t \;\longleftrightarrow\; (t+1)\text{-od-}n \text{ šema}\]
 
\begin{itemize}
\item Tajna $s = f(0)$
\item Udeli su vrednosti polinoma: $s_j = f(j)$
\item Vektor $(s, s_1, \ldots, s_n)$ je kodna reč Rid-Solomonovog koda
\end{itemize}
\end{frame}
 
 \begin{frame}{Algoritam Share}
Neka je $\mathbb{F}_q$ konačno polje, $q > n$, tajna $s \in \mathbb{F}_q$.
 
\vspace{0.8em}
Diler bira nasumičan polinom stepena $t$ sa $f(0) = s$:
\[f(X) = s + f_1 X + f_2 X^2 + \cdots + f_t X^t,
\qquad f_1, \ldots, f_t \overset{\$}{\leftarrow} \mathbb{F}_q\]
 
Učesniku $j$ dodeljuje udeo:
\[s_j \leftarrow f(j), \qquad j = 1, \ldots, n\]
\end{frame}
 
\begin{frame}{Algoritam Recombine}
Skup $Y$ od bar $t+1$ učesnika rekonstruiše tajnu Lagranžovom interpolacijom:
\[s = f(0) = \sum_{i \in Y} s_i \cdot \delta_i(0),
\qquad
\delta_i(0) = \prod_{\substack{j \in Y \\ j \neq i}} \frac{-j}{i-j}\]
 
\begin{itemize}
\item Koeficijenti $\delta_i(0)$ zavise samo od skupa $Y$.
\item Svaki skup sa manje od $t+1$ učesnika ne saznaje ništa o $s$.
\end{itemize}
\end{frame}

 
\begin{frame}{Ilustracija: 3-od-5 šema}
\centering
\begin{tikzpicture}[scale=0.95]
\draw[->] (-0.3,0) -- (6.2,0) node[right]{$x$};
\draw[->] (0,-0.3) -- (0,4.2) node[above]{$f(x)$};
\draw[thick, blue!70!black, domain=0.1:5.5, samples=80]
      plot (\x, {0.22*\x*\x - 0.5*\x + 2.5});
\filldraw[red] (0,2.5) circle (4pt) node[left]{$s = f(0)$};
\foreach \x/\lbl in {1/s_1, 2/s_2, 3/s_3, 4/s_4, 5/s_5}{
\filldraw[blue!70!black] (\x, {0.22*\x*\x - 0.5*\x + 2.5}) circle (3pt)
      node[above right]{\small $\lbl$};
}
\node[align=left, font=\small] at (5.2, 0.7)
     {$t=2$, $n=5$\\potrebno: $\geq 3$};
\end{tikzpicture}
 
\vspace{0.4em}
Bilo koje 3 tačke jednoznačno određuju parabolu, a time i $f(0) = s$.
\end{frame}
 

\begin{frame}{Korekcija pogrešnih udela}
Primenom Berlekamp-Velč algoritma možemo ispraviti pogrešne udele
(učesnici koji lažu).
 
\vspace{0.8em}
Uslov: broj pogrešnih udela $e$ zadovoljava
\[e < t < \frac{|Y|}{3}\]
 
\begin{enumerate}
\item Primeniti Berlekamp-Velč na primljene udele $(j, s_j)_{j \in Y}$.
\item Dobiti ispravljeni polinom $f(X)$.
\item Izračunati $s = f(0)$.
\end{enumerate}
\end{frame}

 
 \section{PRSS}
 
\begin{frame}{Pseudoslučajno deljenje tajni (PRSS)}
\sloppy
\textbf{Cilj:} generisati deljenja novih nasumičnih vrednosti bez interakcije
(osim jednokratne inicijalizacije).
 
\vspace{0.8em}
\textbf{Faza inicijalizacije:} za svaki podskup $A \subset X$, $|A| = n-t$,
distribuira se tajna $r_A$ svim učesnicima iz $A$, zajedno sa
pseudoslučajnom funkcijom $\psi(r_A, \cdot)$.
Definišemo polinom $f_A$ stepena $t$: $f_A(0)=1$, $f_A(i)=0$ za $i \notin A$.
 
\vspace{0.8em}
\textbf{Generisanje deljenja} (bez dilera): biramo javnu vrednost $a$,
učesnik $i$ računa:
\[s_i = \sum_{\substack{i \in A \subset X \\ |A|=n-t}} \psi(r_A, a) \cdot f_A(i)\]
\[s = f(0) = \sum_{\substack{A \subset X \\ |A|=n-t}} \psi(r_A, a)\]
\end{frame}

\begin{frame}{Pitanja}
\begin{enumerate}
\setlength\itemsep{0.6em}
\item Šta je deljenje tajni? Objasniti pojmove:
\begin{itemize}
\item udeo (\textit{share}),
\item kvalifikujući skup,
\item struktura pristupa,
\item pragovna $t$-od-$n$ šema deljenja tajni.
\end{itemize}
 
\item U kompaniji se čuva glavna šifra za pristup kritičnim podacima.
Učesnici su $D$ (direktor), $S$ (sistem administrator),
$B$ (bezbednosni menadžer), $P$ (pravni savetnik).
Važe pravila: $D$ i $S$ zajedno mogu; $D$ i $B$ zajedno mogu;
$S$, $B$ i $P$ zajedno mogu.
Odrediti: (a) minimalne kvalifikujuće skupove,
(b) strukturu pristupa, (c) DNF reprezentaciju.
 
\item Razmatra se Šamirova $(4,7)$ šema deljenja tajni.
(a) Koliki je maksimalni stepen polinoma koji koristi diler?
(b) Koliko učesnika je potrebno za rekonstrukciju tajne?
(c) Da li dva učesnika mogu rekonstruisati tajnu?
(d) Ako se okupe tri učesnika sa ispravnim udelima, kojim
matematičkim postupkom mogu rekonstruisati tajnu?
\end{enumerate}
\end{frame}

\end{document}
